\section{Composition as a Gate}
\label{sec:ch09-composition-as-a-gate}
\index{composition!as a build gate|(}%

Composition is worth almost nothing as something you remember to run. Two
subgraphs stop composing because of a change to one of them, and the person who
made that change is the last person who will think to check the other.

Mosaic's verification script has composed on every pass since
chapter~\ref{ch:the-first-cut-extracting-catalog}, which was the right call a
chapter early. This chapter adds two things to it, and both exist for the same
reason: the chapter prints artefacts, and a printed artefact goes stale
silently.

\subsection*{The composed config is committed}

\index{composition!committing the output}%
\code{federation/supergraph.json} is in the repository at tag \code{ch09}, which
is a deliberate reversal of what
chapter~\ref{ch:how-a-federated-query-actually-runs} did. That chapter's sample
supergraph is a build artefact and is gitignored, on the grounds that it is
derived from two committed schemas and adds nothing.

The difference is that section~\ref{sec:ch09-not-a-schema} prints what is inside
this one. A listing nobody can check out is a listing that has to be taken on
trust, and a book whose whole method is that every listing came from somewhere
does not get to ask for that. So the file is committed and the gate recomposes
and compares:

\begin{minted}[fontsize=\footnotesize]{text}
[ok]   catalog and mosaic compose into one supergraph
[ok]   the composed config matches federation/supergraph.json
\end{minted}

That works only because composition is deterministic, which
section~\ref{sec:ch09-one-command} established by composing twice and comparing.
If \code{wgc} changes the config format in a version bump, that check fails and
says so, which is what should happen to a chapter that documents a format.

\subsection*{The errors are asserted too}

\index{composition!asserting error messages}%
Every composer message in sections~\ref{sec:ch09-satisfiability}
and~\ref{sec:ch09-reading-the-errors} is produced by
\code{scripts/composition-cases.mjs} and checked on every verification run.
Seven cases: the real pair, which composes, and six edits, which do not.

\begin{minted}[fontsize=\footnotesize]{text}
node scripts/composition-cases.mjs --list

baseline            the two schemas as committed, which compose
unsatisfiable-key   mosaic keeps its key but stops answering for it
duplicate-field     both subgraphs declare Product.title, neither says shareable
incompatible-type   a shared field whose return type differs
missing-key         mosaic declares Product and forgets the key
enum-drift          both subgraphs declare ProductCategory with different members
key-mismatch        the two subgraphs key the same entity on different fields
\end{minted}

\index{composition!cases derived from the real schemas}%
Each case is \code{schema/catalog.graphql} and \code{schema/mosaic.graphql} as
committed, with one literal string replacement applied in memory. No fixture
schema exists anywhere in that script, which is the point: an error message
printed from a toy pair of schemas is a fact about the toy.

The edits have a guard on them that I would put in anything of this shape. A
replacement has to match exactly once, and the script fails if it matches zero
times or twice. I added a space inside the \code{@key} on Mosaic's
\code{Product} to see it fire:

\begin{minted}[fontsize=\footnotesize,breaklines]{text}
  FAILED   missing-key         mosaic declares Product and forgets the key
             case "missing-key": the edit to mosaic matched 0 times, expected exactly 1.
             The text it looks for is:
             type Product @key(fields: "id") {
             The committed schema has changed under this case. Fix the case rather than the schema.
\end{minted}

That run reported four failures rather than one, because four of the five cases
edit that same line. The fifth touches a different line and Catalog's schema, so
it went on passing. One character of schema, four stale cases, and every one of
them named.

Without the guard, a future change to \code{Product} would leave those cases
quietly composing the unedited pair and asserting nothing, which is the failure
mode of every test that mutates its own input.

\index{verification!one implementation, two callers}%
It is a node script, and it is the first one in that repository. Both verify
scripts call it rather than each implementing the same mutate, compose and
assert logic in PowerShell and in \code{sh}. That is not a preference. Mosaic
carries two verification scripts that are supposed to check the same things, and
they have already drifted apart once: the \code{sh} one went two chapters
quietly asserting less than its PowerShell twin, which is worse than having no
second gate, because a green run means something different depending on which
one you happened to execute. The cheapest way not to have that problem again is
to have one implementation.

\subsection*{Verify it in Postman}

\index{Postman!checking the composer's input}%
There is nothing in this chapter for Postman to send. \code{wgc router compose}
is a command-line tool reading two files, and no HTTP request anywhere will tell
you whether a graph composes.

What Postman can check is the thing the composer trusts and cannot verify:
whether those two files are still true.

\begin{minted}{graphql}
{ _service { sdl } }
\end{minted}

Send it to port 5101 and to port 5100, exactly as
chapter~\ref{ch:the-first-cut-extracting-catalog} did, and compare each answer
with the committed file that \code{federation/mosaic.yaml} names for that
subgraph. The verification collection does this for both, and the script asserts
it:

\begin{minted}[fontsize=\footnotesize]{text}
[ok]   both subgraphs publish the committed schemas through _service
\end{minted}

\index{composition!what it cannot see}%
That assertion is doing more work than it looks like. The composer's whole view
of your system is two strings. Everything else about a subgraph, whether it is
running, whether its reference resolver works, whether the key it advertises
resolves to anything, whether it answers in under a second, is invisible to
composition and always will be.

So the honest division is this. Composition proves the two schemas can be one
schema. \code{\_service} proves the schemas are the ones the services actually
publish. And the \code{\_entities} calls of
chapter~\ref{ch:the-first-cut-extracting-catalog} prove the subgraphs can answer
for a key, which is the part
chapter~\ref{ch:the-first-cut-extracting-catalog} found composes perfectly while
being completely broken.

Three checks, three different failures, and only the first one is this chapter's.
A build that runs only the first has verified that two files agree.

\medskip
Which is the whole argument for putting composition in the pipeline rather than
in a runbook. It is fast, it needs nothing running, and it catches a class of
breakage that no test of either service can reach. It is also not evidence that
the graph works.
Chapter~\ref{ch:enter-the-router} puts a router in front of this configuration
and asks it a question, which is the first time anything in this book will have
answered one across a boundary without me assembling the answer by hand.

\index{composition!as a build gate|)}%
